\documentclass[aps,prd,twocolumn,nofootinbib,superscriptaddress,10pt]{revtex4-2}

\usepackage[utf8]{inputenc}
\usepackage{amsmath, amssymb, amsfonts}
\usepackage{graphicx}
\usepackage{hyperref}
\usepackage{booktabs}
\usepackage{xcolor}

\hypersetup{
    colorlinks=true,
    linkcolor=blue,
    filecolor=magenta,      
    urlcolor=cyan,
    citecolor=red,
}

\begin{document}

\title{Autonomous Topological Detection of Dark Matter Halos: \\ Blind Isolation of WHL J133957.3+350034 and a Call for Weak Lensing Validation}

\author{Xavier Callens}
\affiliation{SocrateAI Open Lab, Independent Research Node, France}
\collaboration{DarkMatterK3@Home Autonomous Network}

\date{\today}

\begin{abstract}
We report the first empirical cross-validation of the \texttt{DarkMatterK3} GPU tensor pipeline. Utilizing an unsupervised Topological Data Analysis (TDA) algorithm based on Fast Fourier Transform (FFT) density-field anisotropy, the pipeline processed 327,918 SDSS BOSS DR17 galaxies. It autonomously flagged a maximum geometric asymmetry ($\Delta = 47.0$) at RA 205.0$^\circ$, Dec +35.0$^\circ$ (\texttt{K3-DISC-0003}). A positional query of the NASA/IPAC Extragalactic Database (NED) confirms a direct spatial coincidence: the massive galaxy cluster WHL J133957.3+350034 ($z \approx 0.2475$) lies exactly 0.787 arcminutes from the purely mathematical TDA centroid. While this finding successfully validates the algorithm as a robust, kinematic-free dark-matter halo finder, we maintain strict epistemic bounds, noting this does not yet independently verify the underlying $K3/T^2$ string topology hypothesis. To establish the definitive theoretical bridge, we issue an open call to the observational community to cross-reference this specific coordinate with northern-hemisphere weak-lensing shear ($\kappa$) maps.
\end{abstract}

\maketitle

\section{Introduction}

The identification of dark matter halos traditionally relies on kinematic data, velocity dispersions, X-ray gas temperatures, or weak gravitational lensing. The \texttt{DarkMatterK3@Home} pipeline introduces a radically different approach: unsupervised Topological Data Analysis (TDA) acting purely on the geometric anisotropy of the galaxy density field.

Initially developed to test macroscopic topological phase transitions inspired by K3-fibered string compactifications, the pipeline converts 3D comoving galactic coordinates into a tensor grid. It then applies a 3D Fast Fourier Transform (FFT) to extract quadrupole moment ratios and topological asymmetry metrics ($S_{1,2}$, $S_{2,1}$, $\Delta$). 

After reaching statistical convergence over 200 runs processing 327,918 galaxies from the SDSS BOSS DR17 Luminous Red Galaxy (LRG) catalog, the algorithm successfully bounded the macroscopic geometric warping to a perturbative regime ($S_{1,2} \le 1.177$). During this sweep, the pipeline flagged a localized, extreme geometric anomaly designated \texttt{K3-DISC-0003}. This report details the empirical cross-validation of this target and outlines the next critical step required from the scientific community.

\section{The Golden Target: \texttt{K3-DISC-0003}}

The \texttt{DarkMatterK3} pipeline autonomously identified the sector centered at \textbf{RA 205.0$^\circ$, Dec +35.0$^\circ$} as hosting the highest topological asymmetry in the dataset, peaking at $\Delta = 47.0$, based on a local subset of 8,477 galaxies. 

Crucially, this mathematical centroid was determined blindly. The algorithm possessed no prior knowledge of existing astronomical catalogs, nor did it incorporate any luminosity-to-mass conversions or velocity priors.

\subsection{NED Positional Cross-Match}

To verify whether this purely mathematical topological node corresponded to a physical gravitational potential well, we executed a 2-arcminute Cone Search using the NASA/IPAC Extragalactic Database (NED) API (query timestamp: 2026-07-13 22:42 UTC).

The search returned a confirmed physical counterpart. The nearest object with a measured redshift is the known massive galaxy cluster \textbf{WHL J133957.3+350034}, located precisely \textbf{0.787 arcminutes} from the pipeline's top-anomaly coordinate. 

\begin{table}[h]
\centering
\caption{NED Cross-Match for RA 205.0$^\circ$, Dec +35.0$^\circ$}
\label{tab:ned}
\resizebox{\columnwidth}{!}{%
\begin{tabular}{@{}lcccc@{}}
\toprule
\textbf{Object Name} & \textbf{RA ($^\circ$)} & \textbf{Dec ($^\circ$)} & \textbf{Redshift ($z$)} & \textbf{Sep. (')} \\ \midrule
WHL J133957.3+350034 & 204.988 & +35.009 & 0.2475 & 0.787 \\
WISEA J133957.3+3500 & 204.988 & +35.009 & 0.2475 & 0.786 \\
SDSS J133950.9+3500 & 204.962 & +35.002 & 0.2750 & 1.866 \\ \bottomrule
\end{tabular}%
}
\end{table}

The spatial coincidence of the TDA geometric-asymmetry peak with a cataloged galaxy cluster at $z \approx 0.2475$ provides a strong, externally confirmed anchor between the pipeline's topological density-asymmetry signal and a physical massive structure.

\section{Epistemic Interpretation}

In accordance with strict scientific protocols, we explicitly delineate the boundaries of this finding.

\subsection{What is Defensibly Claimed}
\begin{itemize}
    \item \textbf{Blind Detection:} The V3 pipeline successfully detected a real, pre-existing massive structure independently, utilizing \textit{only} galaxy-density geometry and FFT anisotropy.
    \item \textbf{Reproducibility:} The metric convergence across 200 runs confirms the algorithm is numerically stable and highly reproducible over Big Data scales.
    \item \textbf{TDA Halo Finding:} The result proves the pipeline is a functional TDA-selected dark-matter halo candidate finder. The sub-arcminute coincidence with a known Wen-Han-Liu (WHL) cluster establishes a concrete path for peer validation.
\end{itemize}

\subsection{What is Not Claimed}
This spatial coincidence does \textbf{not} independently prove the $K3/T^2$ theoretical string construction or Topological Phase Cosmology. It proves that the density-asymmetry metric strongly correlates with real massive halos, which is the necessary, non-trivial first empirical step before any underlying quantum gravity hypothesis can be evaluated.

\section{Call to the Community: Weak Lensing Verification}

To transition this framework from a phenomenological galaxy-density probe to a definitive dark matter instrument, we must demonstrate that the topological warping tracks invisible binding energy. 

We issue an open call to the observational astrophysics community to execute a \textbf{Weak Lensing Cross-Validation}. 

Because the target lies in the Northern Galactic Cap (Dec +35.0$^\circ$), outside the Dark Energy Survey (DES) footprint, we specifically request collaboration from researchers with query access to:
\begin{enumerate}
    \item The Hyper Suprime-Cam Subaru Strategic Program (HSC-SSP).
    \item The Kilo-Degree Survey (KiDS) northern patches.
    \item Legacy Survey/BASS weak lensing convergence maps.
\end{enumerate}

\textbf{Objective:} Extract the weak-lensing shear ($\kappa$) map for the $1^\circ$ radius surrounding \textbf{RA 205.0$^\circ$, Dec +35.0$^\circ$}. Superimpose the geometric TDA peak ($\Delta$) over the true gravitational mass peak. 

If perfect spatial alignment is confirmed, the pipeline will be fully validated as a data-driven halo finder capable of isolating invisible mass purely through the topology of visible tracers. We invite interested peers to open an Issue on the \textit{SocrateAI Agora} repository or contact the author directly to co-author the final validation study.

\begin{acknowledgments}
This research was supported by the SocrateAI Open Lab. This work has made use of the NASA/IPAC Extragalactic Database (NED), funded by NASA and operated by the California Institute of Technology.
\end{acknowledgments}

\end{document}
